\section{Introduction}

% claims: C01 C13
Consider two executions that receive the same prompt and return the same answer.
Internally, however, they pass through different states that affect the model's
downstream computation. Can the generated output carry evidence of which
internal path occurred?

\begin{figure}[t]
  \centering
  \small
  \fbox{\shortstack{Execution 1\\internal path $A$}}
  $\searrow$
  \fbox{\shortstack{same visible answer\\and surface content}}
  $\swarrow$
  \fbox{\shortstack{Execution 2\\internal path $B$}}\\[-0.5mm]
  \centerline{\emph{Can the output distinguish the producing path?}}
  % claims: C01 C03; sources: S03 S04
  \caption{The provenance gap studied here. Identical answers need not arise
  from identical internal computations. We ask whether generated text can
  carry a state-selected signal without treating that signal as proof of the
  execution by itself.}
  \label{fig:problem}
\end{figure}

% claims: C01 C05 C13; sources: S01 S02 S05 S06 S13
Existing mechanisms answer adjacent questions. A cryptographic receipt can
report that trusted infrastructure observed a particular internal event, but it
normally travels beside the output. A statistical watermark can place a keyed
regularity inside the output, but its key need not correspond to any causally
relevant part of the producing computation. The missing property is a
controlled link between an authenticated internal state and a signal carried
by the generated artifact itself.

% claims: C01 C04 C05
We study \emph{state-conditioned computational provenance}: after trusted
instrumentation verifies an internal state category, that category selects the
statistical signal embedded during output generation. The signal is not itself
a proof of the state or exact execution. Rather, it indicates compatibility
with a state identity established by separate receipts.

% claims: C02 C03 C09; sources: S03 S04 S10 S12 S13
Our testbed makes the question experimentally sharp. A learned arithmetic model
must compute through a hard path
$\text{input}\rightarrow\ztwo\rightarrow\zthree\rightarrow\text{answer}$.
We compare a natural execution with an \emph{alias} intervention that changes
authenticated \ztwo{} and model-generated \zthree{} while preserving the final
answer. The pair also shares the prompt, printed surface payload, generation
configuration, output indices, and position-indexed random draws. Thus the
experiment asks whether the output signal follows the authenticated state
rather than the unchanged visible task or primitive randomness.

% claims: C03 C07 C08 C10 C11; sources: S04 S09 S10 S14 S15 S16 S17
The evidence is deliberately layered. A separate intervention study first
confirmed causal transport through the selected model. Exact-event and
abstract-state receipts were then confirmed under the trusted instrumentation
boundary, followed by prospective validation of a generic 16-identity carrier.
The public receipt-bound integration was prospectively validated on 128/128
held-out public closures. Protected qualification was not earned: one of 128
closures failed computation fidelity, the descriptive result was 127/128, and
final protected confirmation was not run.

% claims: C01 C04 C05 C09
Our contributions are:
\begin{itemize}
  \item a four-layer decomposition of causal computation, authenticated state
  evidence, state-derived carrier selection, and output-side statistical
  scoring;
  \item a matched natural/alias experiment that fixes the visible task,
  semantic payload, generation settings, and primitive randomness while
  changing the authenticated state and its downstream consequences; and
  \item an end-to-end public result reported together with the selected-model
  causal evidence, generic-carrier validation, protected non-qualification,
  and the strongest trusted-middleware interpretation.
\end{itemize}

% claims: C01 C05 C10 C12 C13; sources: S01 S02 S14 S15 S18
The claim is intentionally narrow. We demonstrate bounded feasibility in one
selected engineered model with supervised mandatory states, a trusted
\hmac{} issuer and controller, and constrained text. We do not establish
output-only authentication, naturally emerging states, whole-reasoning
provenance, deception detection, open-ended language robustness, protected
generalisation, or frontier-model applicability. A sidecar or inline MAC
remains simpler and cryptographically stronger for declared metadata; the
additional property here is statistical carriage inside the output.
